% Generated draft: requires mathematical and source review.
\documentclass[11pt]{article}
\usepackage[margin=25mm]{geometry}
\usepackage{fontspec}
\setmainfont{Latin Modern Roman}
\usepackage{amsmath,amssymb,mathtools,unicode-math}
\setmathfont{Latin Modern Math}
\usepackage{enumitem,needspace,xurl,hyperref}
\setlength{\parindent}{0pt}
\setlength{\parskip}{5pt}
\setlength{\emergencystretch}{3em}
\newcommand{\N}{\mathbb{N}}
\newcommand{\Z}{\mathbb{Z}}
\newcommand{\Q}{\mathbb{Q}}
\newcommand{\R}{\mathbb{R}}
\newcommand{\C}{\mathbb{C}}
\newcommand{\F}{\mathbb{F}}
\newcommand{\A}{\mathbb{A}}
\newcommand{\PP}{\mathbb{P}}
\newcommand{\E}{\mathbb{E}}
\newcommand{\eps}{\varepsilon}
\newcommand{\dd}{\,\mathrm{d}}
\begin{document}
\section*{UnsolvedMath ID: 800269}
\textbf{Source problem code:} catalogue entry
\subsection{Definitions and mathematical statement}
For an integer \(k\geq 1\), write \(\PP^k(\C)\) for complex projective \(k\)-space. For a holomorphic self-map \(F:\PP^k(\C)\to\PP^k(\C)\), interpret its degree as its algebraic degree: the common degree \(d\) of homogeneous coordinate polynomials representing \(F\), with no common zero in \(\C^{k+1}\setminus\{0\}\). A point \(p\) is periodic of least period \(n\geq 1\) if \(F^n(p)=p\) and \(F^j(p)\neq p\) for \(1\leq j<n\). It is repelling if every eigenvalue \(\lambda\) of the complex-linear derivative
\[
D(F^n)_p:T_p\PP^k(\C)\longrightarrow T_p\PP^k(\C)
\]
satisfies \(|\lambda|>1\).

Under these conventional interpretations, the first question is whether
\[
\begin{gathered}
\text{for every integer }k\geq 1\text{ and every holomorphic }F:\PP^k(\C)\to\PP^k(\C)\text{ with }\deg F\geq 2,\\
\text{there exist }n\geq 1\text{ and }p\in\PP^k(\C)\text{ of least period }n\text{ such that}
\quad
\operatorname{Spec}\bigl(D(F^n)_p\bigr)\subseteq\{\lambda\in\C:|\lambda|>1\}.
\end{gathered}
\]
The source does not explicitly specify the range of \(k\), the base field, or the convention for degree; the formulation above makes the usual complex-projective interpretation explicit.

The source separately asks whether meromorphic maps have periodic orbits. If this means meromorphic self-maps \(G:\PP^k(\C)\dashrightarrow\PP^k(\C)\), an ordinary periodic orbit would consist of distinct points \(p_0,\ldots,p_{n-1}\), for some \(n\geq 1\), satisfying
\[
p_j\in\PP^k(\C)\setminus I(G),\qquad
G(p_j)=p_{(j+1)\bmod n}\quad(0\leq j<n),
\]
where \(I(G)\) is the indeterminacy locus. This supplies a precise meaning of an orbit avoiding indeterminacy, but the intended class of meromorphic maps remains unspecified: the source does not state explicitly whether the same domain, degree restriction, or a dominance hypothesis is intended, nor whether another convention involving indeterminacy is allowed. Consequently, the second question cannot be given an unqualified universal formulation from the supplied record.
\subsection{Short English statement}
The first question asks whether every holomorphic self-map of complex projective space of algebraic degree at least two has a periodic point at which all eigenvalues of the return-map derivative have modulus greater than one. The source also asks whether meromorphic maps have periodic orbits, but leaves their admissible class and the treatment of indeterminacy unspecified.
\subsection{Sources}
\begin{itemize}
\item Local UnsolvedMath registry, permanent ID \(800269\), source problem code \texttt{catalogue entry}, title: \emph{Several Complex Variables Problem List --- Problem 68}. The supplied registry labels the record open; this label is not an independently verified mathematical or literature assessment.
\item Buzzard, Fornaess, Krantz, \emph{The Several Complex Variables Problem List} (1995), Problem 68. Supplied source URL: \url{https://arxiv.org/abs/math/9509202}. The registry attributes the record to Buzzard, Fornaess, Krantz; the quoted problem ends with the signature \texttt{-sibony@anh.matups.fr}.
\item Supplied provenance records access on \texttt{2026-10-16} and extraction method \texttt{arxiv\_latex}. These are supplied provenance claims; no independent source access or verification is asserted here.
\end{itemize}
\end{document}
